\chapter{异构多核与跨核协同}
\label{chap:heterogeneous-multicore}

现代嵌入式 SoC 常把运行 Linux 的应用处理器与一个或多个实时核、DSP、NPU 或安全岛集成在同一芯片中。异构系统的目标不是简单地“多放一个核”，而是把复杂应用、确定性控制、低功耗常驻和故障隔离分配到最合适的执行域。

\section{任务划分原则}

Linux 适合网络、文件系统、图形界面、复杂协议和动态应用；FreeRTOS 或裸机核适合微秒级控制、持续采样、精确时序和主核休眠期间的常驻任务。划分时应基于截止期、故障影响、内存需求、驱动归属和升级边界，而不是仅按代码语言或团队分工切割。

以下功能通常不宜跨核拆成高频细粒度调用：闭环控制的每个计算步骤、单次寄存器读写和需要共享大量可变状态的算法。跨核接口应尽量是批量、异步、可超时的消息。

\section{执行域契约与分区预算}

每个执行域都应具有可审查契约，而不只是“Linux 核”和“实时核”两个工程目录。契约至少记录：
\begin{itemize}
  \item 职责、输入输出、截止期、数据新鲜度和允许降级；
  \item CPU、SRAM/DDR、共享内存、带宽、功耗和热预算；
  \item 拥有的外设、IRQ、DMA、时钟、复位与电源域；
  \item 启动、暂停、停止、崩溃、升级和回滚时的行为；
  \item 信任级别、可访问地址、密钥和安全输出责任；
  \item 版本兼容窗口、诊断资产和验证方法。
\end{itemize}

推荐把控制平面与数据平面分开。控制平面交换能力、配置、状态和故障；数据平面使用有界批量缓冲区传送采样、帧或推理结果。若所有信息都经过一个无优先级队列，大数据突发可能阻塞停机、心跳和故障确认。

跨域调用的成本不仅是一次中断和 memcpy，还包括排队、cache maintenance、调度、唤醒和上下文恢复。分区预算应使用端到端最大值或高置信上界，并包含对端暂时不可调度、队列接近满和 DDR 竞争条件。平均往返时间不足以决定控制环是否适合跨核拆分。

执行域也是故障与信任边界。Linux 被攻陷时不应自动获得任意控制执行器、重映射远端固件或读取安全岛密钥的能力；实时核崩溃时也不应通过失控 DMA 破坏 Linux 内存。IOMMU、总线防火墙、MPU/PMP、权限受限 mailbox 和最小协议共同形成隔离，单靠“两个团队不会写错”不构成边界。

\section{启动与生命周期}

常见启动模型有：
\begin{itemize}
  \item Bootloader 同时加载 Linux 与实时核固件；
  \item Linux 通过 remoteproc 加载、启动和停止远端核；
  \item 安全固件拥有远端核生命周期，Linux 只能连接通信通道；
  \item 实时核先启动并维持安全输出，Linux 随后接管高级功能。
\end{itemize}

Linux remoteproc 框架负责固件加载、资源声明、启动停止和崩溃报告；RPMsg/OpenAMP 在共享内存 virtio ring 之上提供消息通道\cite{openamp-doc,linux-remoteproc}。资源表通常描述 vring、trace buffer 和设备地址，必须与链接脚本、设备树保留内存和 IOMMU 配置一致。

\begin{center}
\begin{tikzpicture}[node distance=7mm]
  \tikzset{box/.style={draw=bookline,fill=bookpaper,rounded corners=1mm,
    minimum width=38mm,minimum height=10mm,align=center,font=\small}}
  \node[box] (linux) {Linux\\remoteproc/RPMsg};
  \node[box,right=of linux] (shared) {共享内存\\vring＋buffer pool};
  \node[box,right=of shared] (rtos) {FreeRTOS\\OpenAMP endpoint};
  \node[box,below=of shared] (hw) {mailbox/IPI＋硬件信号量};
  \draw[<->,bookteal] (linux)--(shared);
  \draw[<->,bookteal] (shared)--(rtos);
  \draw[<->,bookteal] (linux)--(hw);
  \draw[<->,bookteal] (rtos)--(hw);
\end{tikzpicture}
\end{center}

远端核停止并不等于共享资源已经安全释放。主核必须先阻止新请求，等待在途 DMA 和消息结束，再撤销设备、时钟、内存和中断资源。远端核崩溃后也不能盲目重启，否则可能重复执行一次不可幂等的控制命令。

\subsection{启动 epoch 与能力握手}

每次执行域启动都应生成新的 boot epoch。消息、缓冲区所有权、事务和心跳携带 epoch 或 generation；对端观察到 epoch 改变后，丢弃旧响应、释放陈旧租约并重新协商能力。仅清空 ring 的 head/tail 不足以处理在途 mailbox、中断和 DMA completion。

握手至少交换协议版本范围、feature bits、最大消息/对齐、缓冲区数量、端序、时钟能力、启动原因和当前安全状态。能力协商成功后才开放业务 endpoint；如果核心安全功能不兼容，应保持输出安全并给出稳定错误，而不是降级到一个未验证旧协议。

remoteproc resource table 和远端 ELF 属于解析输入。主核应验证条目数量、长度、对齐、地址范围、vring 大小和 carveout 是否落在允许区域；固件真实性还应由启动链或加载服务验证。签名有效只证明来源和完整性，不证明资源声明不会与 Linux、TEE 或其他远端核重叠。

启动状态机应区分 off、loading、booting、handshaking、ready、quiescing、crashed 和 recovering。业务层不能只通过 \texttt{/dev/rpmsg*} 是否存在判断远端可用，因为设备节点可能早于服务初始化出现，也可能在远端已经失去 progress 时仍然保留。

\section{共享内存与缓存一致性}

共享内存区域必须回答四个问题：双方看到的地址、缓存属性、所有权转换和内存顺序。如果两侧不具备硬件缓存一致性，应在发布描述符前清理数据缓存，在消费前失效缓存，并确保描述符与数据操作顺序正确。

共享布局本身也是持久于一次进程生命周期之外的 ABI。建议在区域头部保存 magic、layout version、总长度、双方 epoch、cache-line 大小、各子区 offset/length 和 header 校验。所有地址使用相对 offset 或双方明确转换的 device address，禁止把任一执行域的虚拟指针写入共享区。

\begin{lstlisting}[language=C,caption={共享内存布局头的逻辑字段}]
struct shared_region_header {
    uint32_t magic;
    uint16_t layout_version;
    uint16_t header_length;
    uint32_t total_length;
    uint32_t host_epoch;
    uint32_t remote_epoch;
    uint32_t vring_offset[2];
    uint32_t pool_offset;
    uint32_t pool_length;
    uint32_t layout_crc;
};
\end{lstlisting}

初始化方先写完整 header 和子区，再执行平台要求的 cache clean/屏障，最后发布 ready generation。\texttt{layout\_crc} 只覆盖不可变布局字段，epoch 通过独立原子更新。连接方只有在验证版本、长度、offset 不重叠、对齐和 epoch 后才能访问子区。共享区来自旧固件或 brownout 后内容未知时，不能因为 magic 恰好匹配就继续使用。

即使硬件一致，编译器和 CPU 仍可能重排普通访问。无锁 ring 需要 release/acquire 语义：生产者先写数据，再以 release 发布写指针；消费者以 acquire 读取写指针后才能访问数据。

\begin{lstlisting}[language=C,caption={共享 ring 的发布与消费语义（伪代码）}]
void ring_publish(struct ring *r, const struct message *m)
{
    unsigned int head = platform_shared_load_relaxed(&r->head);
    r->slot[head % RING_SIZE] = *m;
    platform_cache_clean(&r->slot[head % RING_SIZE], sizeof(*m));
    platform_shared_store_release(&r->head, head + 1);
    platform_notify_peer();
}

bool ring_consume(struct ring *r, struct message *m)
{
    unsigned int head = platform_shared_load_acquire(&r->head);
    if (r->tail == head)
        return false;
    platform_cache_invalidate(&r->slot[r->tail % RING_SIZE], sizeof(*m));
    *m = r->slot[r->tail++ % RING_SIZE];
    return true;
}
\end{lstlisting}

这里的 \texttt{platform\_shared\_*} 不是普通语言原子的别名：在非一致平台上，它还必须对共享 head 所在 cache line 执行正确 clean/invalidate，并使用双方架构都认可的屏障和原子宽度。只清理 payload 而没有发布 head，或只失效 payload 而使用陈旧 head，都会产生静默丢包。

缓存维护函数必须按 cache line 对齐并覆盖完整范围。描述符、head/tail 和高频统计最好放在不同 cache line，避免双方维护同一 line 时把对方字段覆盖回旧值。把共享区设为不可缓存可以简化初期 bring-up，但会显著降低吞吐量，也不能消除对内存屏障的需要。

\subsection{缓冲区所有权与零拷贝}

零拷贝把 memcpy 成本转化为更严格的生命周期契约。每个 buffer 至少经历 free、producer-owned、published、consumer-owned 和 returned 状态；所有权只由单一原子描述符或队列转移，不能一方按超时直接复用而另一方仍在 DMA。

描述符应包含 buffer ID、generation、有效长度、数据类型、时间戳和完整性字段。消费者返回的是精确的 buffer generation，而不是裸索引；这样生产者可以拒绝远端重启前的陈旧归还。发生超时后，系统可标记 buffer 为 quarantined，直到相关 DMA channel、设备和对端 epoch 已确认重置。

对 cache 非一致设备，所有权交接点同时是 cache maintenance 点。生产者发布前 clean，消费者取得后 invalidate；消费者修改并归还时需要反向操作。若某 buffer 同时被 CPU、DSP 和外设访问，应由平台 DMA API 或集中 allocator 建模，禁止三个参与者各自假定自己拥有最新副本。

\section{跨核协议 ABI}

跨核消息是独立固件之间的二进制 ABI，应使用固定宽度字段、显式字节序、版本、长度和事务标识。禁止直接发送 C 指针、编译器相关位域、未初始化 padding 或包含宿主尺寸类型的结构。

\begin{lstlisting}[language=C,caption={可演进的跨核消息头}]
struct ipc_header {
    uint16_t abi_version;
    uint16_t message_type;
    uint32_t total_length;
    uint32_t transaction_id;
    uint32_t source_epoch;
    uint32_t flags;
    uint32_t payload_crc;
};

_Static_assert(sizeof(struct ipc_header) == 24, "IPC ABI drift");
\end{lstlisting}

请求应具有超时、取消和重复响应处理。对可能重试的命令，应定义幂等键或事务号；对事件流，应定义丢包、乱序、背压和消费者重启后的重新同步策略。

传输 endpoint 名称不应直接等于产品服务 ABI。RPMsg 负责发现和传送消息\cite{linux-rpmsg-doc}，服务层仍需定义 authentication/authorization、请求/响应、单向事件、错误码和生命周期。endpoint 重建后，服务必须重新握手，不能沿用旧连接中的权限、订阅或事务状态。

底层 MTU 可能小于业务对象。分片协议应有总长度上限、分片序号、超时、内存预算和重复处理，且在分配大缓冲前先验证元数据。更稳健的做法通常是让消息只描述共享 buffer ID 和有效范围，大载荷走受控 pool；但这又要求严格的所有权与 IOMMU 边界。

能力演进应采用“理解后启用”。发送方只使用双方协商的 feature，未知 message type 可按协议拒绝或忽略，但危险控制命令不能静默降级。字段删除、单位变化和默认行为改变需要新版本或新 type，不能只复用原编号并依赖两侧同步发布。

\texttt{payload\_crc} 只能发现部分随机传输或内存错误，不提供来源认证和抗恶意篡改。若执行域之间存在不同信任等级，应由受保护通道、消息认证码或高权限代理建立身份与授权，并把密钥留在适合的信任边界内。

\section{资源所有权}

同一 UART、I2C、GPIO bank 或 DMA channel 不应被 Linux 和实时核同时驱动。所有权可以静态分配，也可以通过受控状态机迁移，但迁移必须覆盖 pinctrl、时钟、复位、IRQ、DMA 和外设内部状态。

设备树中的 \texttt{reserved-memory} 防止 Linux 页分配器使用远端固件、vring 和共享 buffer；IOMMU 或总线防火墙进一步限制各执行域可访问的地址。硬件信号量只能解决互斥，不能自动解决缓存、超时和崩溃后的锁恢复。

动态迁移建议使用 quiesce---drain---snapshot---reset---reconfigure---grant---resume 状态机：
\begin{enumerate}
  \item 当前所有者拒绝新请求并等待在途事务/DMA 有界结束；
  \item 保存双方真正需要的设备状态，而不是复制全部寄存器；
  \item 屏蔽 IRQ、停止 DMA，并把外设置于无害状态；
  \item 新所有者重新配置 pinctrl、clock、reset、IOMMU 和驱动状态；
  \item 发布新的 ownership generation，旧 generation 的访问全部拒绝；
  \item 完成 loopback/健康检查后恢复业务。
\end{enumerate}

迁移失败时必须有唯一责任方。不能出现旧所有者认为已经释放、新所有者认为尚未获得而外设保持危险输出，也不能双方都因超时重新初始化。硬件 semaphore 可以作为最后防线，但所有权记录、generation 和恢复策略仍由协议定义。

\section{时间同步与时间戳}

双方独立的 tick 计数不能直接比较。需要关联事件时，应选择公共硬件计数器，或通过周期同步估计偏移与漂移。时间同步消息应同时记录发送、接收和处理时刻，避免把通信延迟误认为时钟偏移。

控制闭环最好完全位于实时核；Linux 下发目标值、约束和配置，实时核返回带硬件时间戳的批量采样与状态。这样 Linux 调度抖动不会进入最关键的控制路径。

公共 counter 也需要定义位宽、频率、复位行为和 wrap 处理。若远端在 suspend 中 counter 停止，恢复后必须产生新 time generation；旧 offset 不能继续使用。时间映射记录应包含估计时刻、offset、频率比、最大残差和有效窗口，使消费者能够判断某个时间戳是否仍可比较。

同步协议要抵抗不对称排队。Linux 侧一次高延迟响应可能把 mailbox/调度抖动误估为时钟漂移；可以保存往返时间、丢弃高延迟样本并用多点拟合，但必须给出误差上界。安全或闭环控制若无法接受该不确定度，应使用硬件 latch、同步脉冲或直接共享计数器。

\section{故障检测与恢复}

跨核健康管理至少包括心跳、请求超时、远端 panic trace、看门狗和启动次数。恢复策略应按故障影响区分：重建 endpoint、重启远端核、复位共享外设、重启 Linux 或进入整机安全状态。

远端核更新必须与 Linux 驱动和 IPC ABI 兼容。若两个固件不能独立回滚，应作为一个原子更新集合签名和部署；若允许独立升级，则必须提供能力协商和至少一个兼容窗口。

心跳必须与 progress 区分。远端调度器可以持续发送心跳，但控制任务已经死锁或采样序列停止。健康消息应报告关键 pipeline 的 generation、最后完成事务、队列水位、错误锁存和本地 watchdog 状态；Linux 只在这些指标满足产品门槛时认为远端 ready。

远端 crash 后，主核先冻结新业务并保存 panic/trace、resource table、epoch、队列和设备状态，再决定恢复层级。若直接清零共享区和重启，可能丢失根因并让旧 DMA 写入新固件内存。恢复前应 reset/隔离能够成为 bus master 的远端核和外设，并使输出进入危害分析规定状态。

请求超时不表示对端没有执行。对执行器、计费或一次性写入，重试必须使用幂等 transaction ID，并允许查询最终状态。远端重启后若无法证明旧命令是否完成，应选择保守业务策略，而不是无条件重放。

\section{低功耗、暂停与唤醒协同}

Linux suspend、远端常驻和 SoC 电源域之间需要独立契约。进入低功耗前应协商哪些任务继续运行、共享 DDR 是否保持、mailbox/IRQ 是否可唤醒、公共时钟是否停止，以及唤醒后由谁恢复外设。

若远端在 Linux suspend 期间写共享 DDR，必须保证该内存、电源和 interconnect 仍可用；否则应切换到 retention SRAM 或本地缓冲。Linux 恢复后先验证远端 epoch 和数据时间戳，再消费期间积累的数据，不能假定 ring 与 suspend 前连续。

唤醒原因需要端到端确认。远端产生 wake request、PMIC/SoC 接受、Linux resume 和业务服务 ready 是不同阶段；每段应有超时和错误码。频繁虚假唤醒会破坏能耗预算，也可能表示线路抖动、未清中断或协议 generation 错误。

远端固件更新、调试或崩溃恢复期间应阻止进入不兼容的低功耗状态。电源管理器不能只看 endpoint 在线，而应消费明确的 lifecycle state 和最坏恢复时间。

\section{安全边界与不可信输入}

跨核共享内存和 mailbox 应按不可信输入解析，即使两个固件属于同一产品。长度、offset、枚举、buffer ID、feature 和事务号都要验证；任何来自对端的地址必须限定在已授权 carveout/IOMMU mapping 内。协议解析错误不应让高权限核任意读写物理内存。

Linux 用户空间若能通过 rpmsg char 访问 endpoint，还需由内核驱动或代理服务实施设备节点权限、消息类型白名单、速率和资源配额。不能因为通信“不经过网络”就忽略认证授权；被攻陷的普通应用可能利用远端高权限控制外设。

远端加载接口应限制允许固件、签名、版本和目标核，阻止非授权用户停止安全岛或替换固件。debug trace 也可能泄漏密钥、传感数据和物理地址，应有生命周期权限、容量限制和清理策略。

故障隔离需要验证 DMA 防火墙，而不只是检查配置寄存器。测试可让远端尝试访问 carveout 边界外地址，确认 IOMMU fault、总线错误、日志和恢复均符合预期，同时保证注入不会破坏其他执行域。

\section{跨域可观测性与更新}

日志应携带执行域 ID、boot epoch、单调 sequence、公共/本地时间戳和 transaction ID。文本时间相同不能证明事件顺序；主核收集远端 trace 时还要记录采集延迟、丢包和时间映射版本。统一事件语义可结合第~\ref{chap:observability-fleet}~章的方法。

发布 manifest 应把 Linux、Device Tree、remote firmware、resource-table/layout ABI 和服务协议版本绑定。若允许独立更新，兼容矩阵至少覆盖新 Linux/旧远端、旧 Linux/新远端、回滚和数据 schema；若不能支持这些组合，就应作为同一更新事务安装和回退。

新远端固件的健康确认不能只看 endpoint 出现。应验证 capability、共享内存压力、外设 loopback、控制输出安全和规定时间内 progress；确认前保留旧固件或整包回滚能力。防回滚计数还需与整组兼容策略一致，避免只升级远端后使整机无法回退到可工作的组合。

\section{Linux/FreeRTOS 协同实现骨架}

一个最小可运行链路包括四项：Device Tree 保留共享内存并声明 remoteproc；远端 ELF 携带 resource table；FreeRTOS 建立 OpenAMP endpoint；Linux RPMsg 驱动或用户空间 rpmsg char 发送消息。

远端资源表声明两个 vring 和 trace buffer，地址既可以由 Linux 动态分配，也可以固定到双方约定的 carveout：

\begin{lstlisting}[language=C,caption={远端固件 resource table 结构（简化）}]
struct remote_resource_table {
    struct resource_table base;
    uint32_t offsets[2];
    struct fw_rsc_vdev rpmsg_vdev;
    struct fw_rsc_vdev_vring vring0;
    struct fw_rsc_vdev_vring vring1;
    struct fw_rsc_trace trace;
};

__attribute__((section(".resource_table")))
const struct remote_resource_table resources = {
    .base = { .ver = 1, .num = 2 },
    .offsets = {
        offsetof(struct remote_resource_table, rpmsg_vdev),
        offsetof(struct remote_resource_table, trace),
    },
    .rpmsg_vdev = { .type = RSC_VDEV, .id = VIRTIO_ID_RPMSG,
                    .num_of_vrings = 2 },
    .vring0 = { .da = FW_RSC_ADDR_ANY, .align = 0x1000, .num = 256 },
    .vring1 = { .da = FW_RSC_ADDR_ANY, .align = 0x1000, .num = 256 },
    .trace = { .type = RSC_TRACE, .da = TRACE_BUFFER_DA,
               .len = TRACE_BUFFER_SIZE, .name = "rtos-trace" },
};
\end{lstlisting}

FreeRTOS 初始化平台 mailbox、共享内存和 libmetal 后创建 RPMsg endpoint。回调只验证消息并投递到任务队列，不在 OpenAMP 回调上下文执行长事务：

\begin{lstlisting}[language=C,caption={FreeRTOS 端 RPMsg 接收边界}]
static int endpoint_rx(struct rpmsg_endpoint *ept, void *data,
                       size_t len, uint32_t src, void *priv)
{
    struct ipc_message msg;

    if (len < sizeof(struct ipc_header) || len > sizeof(msg))
        return RPMSG_SUCCESS;
    memcpy(&msg, data, len);
    if (!ipc_validate(&msg, len))
        return RPMSG_SUCCESS;
    if (xQueueSend(ipc_queue, &msg, 0) != pdPASS)
        atomic_fetch_add(&ipc_dropped, 1);
    return RPMSG_SUCCESS;
}
\end{lstlisting}

Linux 启动远端核后，应等待 endpoint 出现并执行版本握手，再开放业务设备：

\begin{lstlisting}[language=bash,caption={Linux 侧启动与握手检查（示意）}]
echo realtime-fw.elf > /sys/class/remoteproc/remoteproc0/firmware
echo start > /sys/class/remoteproc/remoteproc0/state
wait-for-device /dev/rpmsg_ctrl0 --timeout 5
ipc-selftest --endpoint control --require-abi 3
\end{lstlisting}

测试必须覆盖 Linux 先启动、实时核先启动、endpoint 重建、队列满、消息 CRC 错误、远端 panic、Linux 重启和两侧版本不匹配。共享外设在远端恢复前应保持安全状态，不能仅凭 endpoint 再次出现就恢复输出。

实现还应提供 host-only 协议测试和目标端共享内存自检。host-only 测试可穷举长度、版本、分片、重复事务和状态机；目标端自检则验证实际地址转换、cache line、barrier、mailbox、IOMMU 和 DMA。两者复用同一 schema 与 golden vector，避免模拟器和固件分别解释协议。

\section{验证清单}

\begin{itemize}
  \item 每个执行域是否具有职责、资源、信任、生命周期和升级契约？
  \item 冷启动、远端晚启动、重复启动和停止是否都具有确定结果？
  \item 重启 epoch 改变后，陈旧消息、buffer、DMA completion 和事务是否被拒绝？
  \item 共享内存在一致和非一致平台上的 cache 操作是否经过压力验证？
  \item 共享布局的长度、offset、对齐、cache line 和所有权 generation 是否机器校验？
  \item 消息队列满、远端无响应、重复响应和主核重启是否能够恢复？
  \item 外设、DMA、IRQ 和时钟是否始终只有一个明确所有者？
  \item 动态所有权迁移是否覆盖 quiesce、在途 DMA、失败回退和安全输出？
  \item 两侧版本不匹配时是否拒绝危险操作并输出可诊断原因？
  \item Linux suspend/resume 与远端常驻任务之间是否有明确契约？
  \item IOMMU/总线防火墙是否用越界 DMA 故障注入验证，而非只检查配置？
  \item 心跳是否同时证明关键 pipeline progress，而非仅证明调度器存活？
  \item Linux、DTB、远端固件、layout 与服务 ABI 是否具有可回滚兼容矩阵？
  \item 远端核崩溃时，执行器是否进入系统危害分析定义的状态？
\end{itemize}

第~\ref{chap:realtime-language-memory-model}~章进一步解释 release/acquire、barrier、DMA 所有权和实时阻塞边界；第~\ref{chap:system-build-update}~章给出整组镜像兼容与原子更新方法。
